\chapter{Methodology}
\label{ch:methodology}

This chapter details the systematic procedures, processes, and tools utilized in the design, development, and evaluation of the ScholarPath AdDU system. It outlines the research design, system development methodology, system architecture, data gathering instruments, and statistical tools used for data analysis.

\section{Research Design}
\label{sec:research-design}

This study employs a Developmental Research Design, specifically focusing on System Development. This approach is the most appropriate for IT and Computer Studies capstone projects, as it centers on the structured creation, implementation, and evaluation of a functional software solution to address a specific operational problem. In this context, the design governs the transition of Ateneo de Davao University's (AdDU) paper-based scholarship admission process, as defined in the Standard Procedure for Scholarship Applications (SOP) \cite{ref51}, into a digital, data-feeding workflow that supports the University Scholarship Office and the five School Scholarship Subcommittees. The system requirements are drawn from Key Informant Interviews (KII) with the admissions office and are consolidated in the project's ISO audit preparation report \cite{ref52}. The developmental design is complemented by a descriptive evaluation phase to assess the prototype's usability and functional suitability.

The research design incorporates a comparative usability evaluation to validate the system's impact on processing efficiency. The researchers will administer the Pre-Study Student Problem Validation Survey (Appendix~\ref{app:pre-survey}) to a purposive sample of current or recent AdDU scholarship applicants. This survey measures the perceived time and effort applicants spent on the paper-based process, from reading the scholarship announcement through preparing and physically submitting the SOP application form and its supporting documents. % TODO(L-2): confirm whether Form 230-SCH is the SOP application form.
Once the ScholarPath AdDU prototype is developed, a controlled, task-based usability testing session will be conducted. The perceived time and effort recorded at baseline will then be compared against the corresponding items in the Post-Prototype Perceived Effort Survey (Appendix~\ref{app:post-survey}), which is administered to the same participants immediately following the prototype session. Both instruments use parallel ordinal and Likert-scaled items, enabling direct item-to-item comparisons to determine whether participants perceive the two-stage digital workflow to require less time and effort than the manual process they described at baseline. University Scholarship Office staff and School Scholarship Subcommittee members who complete the staff-side tasks additionally complete the Committee Portal Usefulness Survey (Appendix~\ref{app:committee-survey}). The design and administration of all instruments are detailed in Section~\ref{sec:instruments}.

The comparative evaluation follows a within-subjects design, meaning the same cohort of applicant participants who complete the Pre-Study Survey subsequently participate in the prototype usability testing session. This design is chosen over a between-subjects approach because it eliminates inter-group variability. Each participant serves as their own control, so any observed difference in perceived effort is attributable to the system rather than to individual differences in prior scholarship experience. A sample of 10 to 15 participants is sufficient for this design given that SUS reliability stabilizes at that range \cite{ref42}, and the time-perception comparison uses ordinal scaled data that does not require large-N inferential statistics.

\section{System Development Life Cycle (SDLC)}
\label{sec:sdlc}

To ensure flexibility and iterative improvement, the development of ScholarPath AdDU follows the Agile Development Methodology. Unlike rigid, linear models, Agile allows the development team to adapt to evolving requirements from stakeholders, including the University Scholarship Office, the Office of Admission, and the School Scholarship Subcommittees, through continuous feedback loops between phases. As illustrated in Figure~\ref{fig:sdlc}, the development process is organized into five iterative phases, each building upon the outputs of the last.

\begin{figure}[htbp]
  \centering
  \includegraphics[width=0.8\textwidth]{figure2_p23.png}
  \caption{ScholarPath AdDU Agile SDLC Flowchart}
  \label{fig:sdlc}
\end{figure}

Phase 1: Requirements Gathering. The first phase involves data collection through Key Informant Interviews (KII) with primary stakeholders, including admissions and University Scholarship Office personnel. This phase identifies the eight steps of the SOP \cite{ref51}, the documents required at each step, the selection criteria applied by each School Scholarship Subcommittee, and the points in the process where a person decides an applicant's status. Phase 2: System Design. Drawing from the gathered requirements, this phase produces the architectural blueprint of the system. Deliverables include the relational database schema, User Interface/User Experience (UI/UX) wireframes for applicant, staff, and committee views, and the application status state machine that governs the Application Lifecycle Path. Phase 3: Development. The core programming phase translates design artifacts into a functioning system. Front-end interfaces are built using Vite as the primary build tool, paired with a component-based JavaScript framework (React or Vue) for a reactive, modular UI. Back-end infrastructure, including the PostgreSQL database, secure Document Vault storage, and user authentication, is managed entirely through Supabase. To extend the application to mobile platforms, Capacitor wraps the Vite production build into native Android and iOS packages, enabling cross-platform deployment without maintaining a separate codebase. Phase 4: Testing. Prior to wider release, the system undergoes alpha testing to verify that core functionalities, such as two-stage document uploading, Conditional Revert handling, and human-confirmed status changes, operate without critical errors. Issues surfaced during this phase are documented and fed back into development for resolution before proceeding. Phase 5: Deployment and Evaluation. The functional prototype is deployed to a purposive sample of end-users for User Acceptance Testing (UAT). System quality is then assessed using standardized usability metrics defined by the ISO/IEC 25010 standard, providing quantitative and qualitative data to inform further refinement.

As the feedback loop in Figure~\ref{fig:sdlc} indicates, findings are not terminal; they inform subsequent iterations. The present design is itself the product of one such Adaptation cycle. An earlier prototype automatically assigned students to grants from academic and income data. Following the KII with the admissions office and preparation for an institutional ISO audit \cite{ref52}, that automated matching engine was removed. It was replaced by a data-feeding architecture in which no module assigns, approves, or disapproves an applicant without a recorded human action.

\section{System Architecture and Algorithmic Logic}
\label{sec:architecture}

The architectural decisions underlying ScholarPath AdDU respond to two sources of evidence: the operational requirements documented in the SOP \cite{ref51} and the KII \cite{ref52}, and the gaps identified in the reviewed literature. The choice of a centralized Supabase/PostgreSQL backend was motivated by the policy desynchronization documented in Catibog \cite{ref16} and Orgianus et al. \cite{ref36}, where decentralized systems produced contradictory information that a single canonical database is architecturally positioned to eliminate. The Document Vault and its two-stage submission design directly address the integrated document handling gap identified by both Al-Ayyubi and Maulana \cite{ref3} and Bagaporo and Espinosa \cite{ref12}, whose systems still required physical document submission as a separate, offline step. Bilog's \cite{ref14} finding that unaided self-assessment of eligibility drove application errors motivates Phase~1 Pre-qualification, in which the system flags baseline criteria early and a University Scholarship Office staff member confirms the outcome. Kanwal et al.'s \cite{ref50} acknowledgment that machine-learned eligibility decisions offer no interpretable audit trail motivates the central governance rule of the platform: the system structures, filters, and presents applicant data, and a person makes every decision that changes an applicant's status. The automated notification subsystem responds to the missed-deadline patterns attributable to passive communication channels documented by Catibog \cite{ref16}. Finally, the hybrid mobile-and-web architecture via the Capacitor wrapper was motivated by the World Bank's \cite{ref48} caution that digitization without mobile-first design fails to achieve adoption in Philippine HEI contexts.

\subsection{Scope: Internal Scholarship Tracks}
\label{sec:tracks}

The platform governs the three internal scholarship tracks through which AdDU selects approximately 1,326 applicants per cycle: the Jubilee Scholarship, the Grant-in-Aid (GIA), and the Working Scholars program (Table~\ref{tab:tracks}, Chapter~\ref{ch:introduction}). At registration, each applicant chooses one internal track and up to two degree programs, matching their entrance exam selections. External and government grants, such as those from the Department of Science and Technology (DOST), local government units, and private foundations, are outside the selection process. The university processes their disbursement through financial obligation clearance only after receiving the official beneficiary roster, and the platform treats them as a non-selection pathway.

\subsection{Application Lifecycle Path and Two-Stage Submission}
\label{sec:lifecycle}

The name ScholarPath refers to two structured paths. The Application Lifecycle Path carries an applicant from registration through Phase~1 Pre-qualification, Phase~2 Full Verification, interview scoring, and subcommittee selection, to an award, waitlist, or disapproval outcome. The Academic Trajectory Path then tracks an awarded scholar's academic standing, the completion of service obligations for Working Scholars, and graduation within the contracted degree program. Figure~\ref{fig:lifecycle} presents the end-to-end Application Lifecycle Path.

\begin{figure}[htbp]
  \centering
  % TODO(FIG): replace this placeholder with
  %   \includegraphics[width=\textwidth]{figure_lifecycle.png}
  % drawn to match the ISO audit report's lifecycle diagram (10 steps, 1 loop back).
  \fbox{\parbox{0.92\textwidth}{\centering\small\vspace{0.8em}
    Registration $\rightarrow$ Phase~1 Pre-qualification $\rightarrow$ Phase~2 Full Verification $\rightarrow$ Interview $\rightarrow$ School Subcommittee Selection $\rightarrow$ Release of Results (Office of Admission) $\rightarrow$ Applicant Notified $\rightarrow$ Academic Trajectory Path\\[0.6em]
    Full Verification $\xrightarrow{\text{missing files}}$ Conditional Revert $\xrightarrow{\text{corrected files}}$ Full Verification\\[0.3em]
    Conditional Revert $\xrightarrow{\text{not corrected after 5 days}}$ Lapsed (staff confirm disapproval)\vspace{0.8em}}}
  \caption{Application Lifecycle Path of ScholarPath AdDU}
  \label{fig:lifecycle}
\end{figure}

The revised workflow keeps all eight SOP steps in their original order. It changes only how documents are submitted (step~2) and adds system support for steps~3 to~8, as summarized in Table~\ref{tab:sop-alignment}. % TODO(OPEN-10): will the SOP be formally revised for two-stage submission and Conditional Revert?

\begin{table}[htbp]
  \centering
  \small
  \caption{Alignment of ScholarPath AdDU with the Standard Procedure for Scholarship Applications}
  \label{tab:sop-alignment}
  \begin{tabular}{@{}p{3.6cm}p{6.2cm}p{4.6cm}@{}}
    \toprule
    SOP Step & System Function & Change from SOP \\
    \midrule
    1. Application announcement & Public information page and inquiry form & None \\
    2. Submission of applications & Two-stage digital submission to the Document Vault & Single submission becomes two phases \\
    3. Document verification & University Scholarship Office reviews files in the Document Vault and encodes applicant profiles & None \\ % TODO(OPEN-8): who performs step 3 in practice.
    4. Endorsement to next stage & Qualified applications are released to interview scheduling & None \\
    5. Interview & Interview panel assigned per application records scores and recommendations & Panel composition is configurable \\ % TODO(OPEN-2)
    6. Evaluation and deliberation & School Scholarship Subcommittee reviews data presented by the system and records its decision & None \\
    7. Recommendation and approval & The subcommittee's accepted list is sent to the Office of Admission & System records the handoff \\
    8. Release of results & Office of Admission releases results; applicants are notified only after that release & Notification follows the release event \\
    \bottomrule
  \end{tabular}
\end{table}

Applicants submit documents in two windows (Table~\ref{tab:phases}). In Phase~1 Pre-qualification, the system screens basic eligibility, and pre-qualified applicants receive a provisional go-ahead without submitting a full physical file. In Phase~2 Full Verification, the University Scholarship Office completes verification of the remaining documents before committee review. All credential verification takes place inside the University Scholarship Office; no external or third-party verification service appears in any workflow.

\begin{table}[htbp]
  \centering
  \small
  \caption{Two-Stage Submission Windows and Required Documents}
  \label{tab:phases}
  \begin{tabular}{@{}p{2.6cm}p{2.2cm}p{9.6cm}@{}}
    \toprule
    Phase & Window & Documents \\
    \midrule
    At registration & --- & Application form \\
    Phase 1: Pre-qualification & January to April & Entrance exam result with an ``ACCEPTED'' remark; high school report card (second and third quarter) with a general average of at least 85\%; household income proof; Certificate of Award as Rank~1 or Rank~2 (Jubilee applicants) \\
    Phase 2: Full Verification & May to enrollment & Proof of residence or photographs of residence; recommendation from the homeroom moderator; Certificate of Good Moral Character; BIR Form 2316 or sworn statement of income (GIA and Working Scholars); additional documents configured per cycle, such as a personal essay or an ethnolinguistic group endorsement where applicable \\ % TODO(OPEN-6): residence photos vs. proof; essay and ethnolinguistic endorsement still required?
    \bottomrule
  \end{tabular}
\end{table}

\subsection{Three-Layer Screening and Committee Governance}
\label{sec:governance}

The system structures, filters, and presents applicant data, and people make every decision that changes an applicant's status. Table~\ref{tab:decision-boundary} sets this boundary for each function of the platform. Every status change is written to an audit log that records the acting user, the timestamp, and the prior and new status, so that each decision can be traced to a person during an institutional audit.

\begin{table}[htbp]
  \centering
  \small
  \caption{Decision Boundary Between System Functions and Human Decision-Makers}
  \label{tab:decision-boundary}
  \begin{tabular}{@{}p{3.4cm}p{6.6cm}p{4.4cm}@{}}
    \toprule
    Function & What the System Does & Who Decides \\
    \midrule
    Eligibility screening & Checks baseline criteria (85\% high school average, entrance exam remark, required documents) and flags results & University Scholarship Office staff confirm \\
    Document verification & Stores uploads and shows a checklist per application & University Scholarship Office staff, internally only \\
    Conditional Revert & Sends the notice, runs the 5-day clock, and marks a file ``Lapsed'' & Staff confirm any disapproval \\
    Interview scoring & Records panel scores and recommendations & Interview panel \\
    Selection & Presents applicant data, sorted and filtered by criteria the committee sets & School Scholarship Subcommittee \\
    Waitlist and slot reallocation & Alerts the committee to a vacancy and proposes the next waitlisted candidate & School Scholarship Subcommittee \\
    Track realignment and program transfer & Records the request and checks slot availability & Evaluators and the receiving school \\
    \bottomrule
  \end{tabular}
\end{table}

Evaluation proceeds through three layers. In Layer~1 (General System Verification), the system checks the baseline criteria: an entrance exam result with an ``ACCEPTED'' remark, a high school general average of at least 85\%, and the required documents. An applicant below the baseline, for example with an 84\% average, is not blocked at registration. The application enters a staff review queue, where a University Scholarship Office staff member records a disapproval or waitlist outcome and the applicant receives a notice.

In Layer~2 (Interview Evaluation and Recommendation), interview panels record scores and submit formal recommendations. Following SOP step~5, the interview panel is assigned per application, and panel composition is a configurable setting of each cycle. % TODO(OPEN-2): school-based panels (SOP) vs. Scholarship Director and designated interviewers.
The committee may configure a shortlisting rule that advances a proportion of the interviewed pool, such as the top 50\%, to Layer~3; the system applies no cut unless the committee enables one. % TODO(OPEN-7): is the top-50% cut a formal rule?

In Layer~3 (School Subcommittee Selection), the system partitions applications by degree choice and presents each School Scholarship Subcommittee's data under that school's criteria (Table~\ref{tab:subcommittees}). An applicant whose two program choices fall under different schools appears in both partitions. The system encodes each school's criteria as a rule set that produces flags and sort orders for the committee; it does not produce a selection.

\begin{table}[htbp]
  \centering
  \small
  \caption{School Scholarship Subcommittee Rule Sets}
  \label{tab:subcommittees}
  \begin{tabular}{@{}p{4.2cm}p{3.6cm}p{6.6cm}@{}}
    \toprule
    Subcommittee & Academic Focus & Selection Criteria Presented by the System \\
    \midrule
    School of Nursing (SON) & Nursing and allied health & Zero-conditional policy: conditional entrance exam scores are flagged for rejection; an optional 1-to-10 scoring matrix is available \\
    School of Engineering and Architecture (SEA) & Engineering and architecture programs & Hard cutoffs on entrance exam scores; applicants below the minimum are flagged as not meeting the program requirement regardless of high school rank \\ % TODO(OPEN-4): which school houses Computer Science and IT.
    School of Business and Governance (SBG) & Business and management & Entrance scores combined with financial need tiers and qualitative leadership and potential assessments \\
    School of Education (SOE) & Teacher education & Academic performance weighed alongside commitment to community teaching service \\
    School of Arts and Sciences (SAS) & Humanities, social sciences, and sciences & Academic potential and alignment with university leadership objectives, coordinated across the school's departments \\ % TODO(OPEN-4): four departments with Assistant Deans?
    \bottomrule
  \end{tabular}
\end{table}

The Jubilee Scholarship grant rate is fixed once a Rank~1 or Rank~2 certificate is verified. Program-specific admission minimums remain configurable per school, and when a Jubilee applicant does not meet a school's hard minimum, the system flags the conflict and the subcommittee's recorded decision stands. % TODO(OPEN-3): does the Jubilee guarantee override a hard entrance exam cutoff? TODO(OPEN-9): is the SOP "academic grant" the Jubilee track?

\subsection{Conditional Revert Protocol}
\label{sec:revert}

A missing or illegible document no longer produces an abrupt rejection. Instead, every incomplete file receives five calendar days to be corrected before a staff member decides its outcome:

\begin{enumerate}
  \item A University Scholarship Office staff member finds a missing or illegible document during verification and sets the application status to Conditional Revert.
  \item The platform sends an automated notice that lists each missing or illegible file.
  \item The applicant has five calendar days to upload corrected files. The clock starts at a configurable event, by default the moment the notice is sent. % TODO(OPEN-8): when does the 5-day clock start?
  \item Corrected files return the application to verification.
  \item If nothing arrives within five days, the system marks the file ``Lapsed'' and alerts staff. A University Scholarship Office staff member confirms the disapproval.
\end{enumerate}

Listing~\ref{lst:lifecycle} specifies the application status state machine that enforces this protocol and the governance rule of Section~\ref{sec:governance}. Automated events may only move an application into the Conditional Revert or Lapsed states or raise alerts. Every transition into an outcome state requires a recorded human actor.

\begin{lstlisting}[caption={Pseudocode of the application status state machine with human-confirmation guards},label={lst:lifecycle}]
FUNCTION ChangeStatus(Application, NewStatus, Actor, Reason):
    // Guard 1: outcome states always require a human actor
    IF NewStatus IN {Disapproved, Waitlisted, Awarded, Endorsed} THEN
        RequireHumanActor(Actor)       // reject if Actor is the system
        RequireRole(Actor, AllowedRoles[Application.Status, NewStatus])
    END IF

    // Guard 2: only permitted transitions are accepted
    IF (Application.Status, NewStatus) NOT IN AllowedTransitions THEN
        REJECT "Transition not permitted"
    END IF

    AuditLog.APPEND(Application.ID, Actor, CURRENT_TIMESTAMP,
                    Application.Status, NewStatus, Reason)
    SET Application.Status = NewStatus
    EmitEvent("StatusChanged", Application)
ENDFUNCTION

// Scheduled daily by the Cron worker (system actor)
FUNCTION CheckRevertClocks():
    FOR EACH Application WHERE Status == ConditionalRevert:
        IF CURRENT_DATE > Application.RevertStart + 5 DAYS
           AND Application.HasPendingCorrections THEN
            SET Application.Status = Lapsed    // flag only, not a decision
            AuditLog.APPEND(Application.ID, "system", CURRENT_TIMESTAMP,
                            ConditionalRevert, Lapsed, "Revert window elapsed")
            AlertStaff(Application)            // staff confirm disapproval
        END IF
    ENDFOR
ENDFUNCTION

// Triggered when an awardee declines
FUNCTION OnAwardDeclined(Application):
    FreeSlot(Application.Program, Application.Track)
    Candidate = NextWaitlisted(Application.School, Application.Track)
    AlertSubcommittee(Application.School, ProposedCandidate = Candidate)
    // No award is made until the subcommittee calls ChangeStatus
ENDFUNCTION
\end{lstlisting}

\subsection{Scholar Movement, Track Realignment, and Slot Management}
\label{sec:movement}

The system supports three movement protocols (Table~\ref{tab:movement}). In each, the system records the change and alerts the responsible party immediately, and a person makes the decision.

\begin{table}[htbp]
  \centering
  \small
  \caption{Scholar Movement Protocols}
  \label{tab:movement}
  \begin{tabular}{@{}p{2.6cm}p{4.3cm}p{4.0cm}p{3.5cm}@{}}
    \toprule
    Protocol & Example & System Action & Decision By \\
    \midrule
    Inter-track realignment & A Working Scholar candidate with high financial need and strong academic merit is moved to Grant-in-Aid & Records the change and keeps the history of the applicant's track & Evaluators during screening \\
    External grant program transfer & A DOST awardee must enroll in a specific STEM field, so the applicant moves from Nursing to an engineering program & Checks slot availability in the destination program and records the transfer & Receiving school, subject to slot availability \\
    Waitlist and slot reallocation & A primary awardee declines after receiving an offer from another institution & Frees the slot, alerts the subcommittee, and proposes the next waitlisted candidate & School Scholarship Subcommittee confirms; the candidate is notified after confirmation \\
    \bottomrule
  \end{tabular}
\end{table}

\subsection{Committee Filtering and Backend Indexing Strategy}
\label{sec:filtering}

To let the University Scholarship Office and each School Scholarship Subcommittee review approximately 1,326 applications per cycle without manual collation, the Data-Feeding Committee Portal exposes applicant records through intersecting filters (e.g., ``School = SEA'' + ``Track = GIA'' + ``Verification Status = Complete''). The available facets are school, track, program choice, lifecycle status, verification status, interview score, and screening flags. The sort order and filter criteria are set by the committee, not by the system. Indexing Strategy for Real-Time Retrieval: to ensure real-time rendering, the Supabase backend uses compound indexing, pre-sorting records on frequently queried combinations instead of scanning the entire table for every query.

\begin{itemize}
  \item B-Tree Indexes are applied to exact-match categorical fields (e.g., \texttt{School\_Code}, \texttt{Track}, and \texttt{Status}).
  \item Range Indexes are applied to continuous numerical and temporal fields (e.g., \texttt{HS\_Average}, \texttt{Interview\_Score}, and \texttt{Revert\_Start}).
  \item Client-Side Caching and Debouncing: The frontend interface uses asynchronous JavaScript (Fetch) combined with an input debounce of 300 milliseconds, so the database is queried only after the user finishes selecting facets, reducing server load and preventing interface stutter.
\end{itemize}

\subsection{System Architecture and Cloud Infrastructure}
\label{sec:cloud}

To resolve the operational bottlenecks caused by the physical submission of dense paperwork, ScholarPath AdDU employs a secure, cloud-based architecture. The integration of the database, storage, and communication services is mapped in Figure~\ref{fig:architecture}.

\begin{figure}[htbp]
  \centering
  % TODO(FIG): redraw figure4_architecture.png -- relabel the client apps as
  % "Applicant/Scholar App" and "Committee Portal"; add an Audit Log store in
  % the backend layer; keep no external verification service.
  \includegraphics[width=0.9\textwidth]{figure4_architecture.png}
  \caption{High-Level System Architecture of ScholarPath AdDU}
  \label{fig:architecture}
\end{figure}

Technical Stack and Storage: The system is built on a three-layer technical stack selected for its cohesion and rapid prototyping capabilities. At the client layer, Vite serves as the frontend build tool, paired with a component-based JavaScript framework (React or Vue), and provides two interfaces: the applicant and scholar application, and the Data-Feeding Committee Portal used by University Scholarship Office staff, interview panels, School Scholarship Subcommittees, and the Office of Admission. For cross-platform mobile deployment, Capacitor wraps the same Vite production build into native Android and iOS applications. At the backend layer, Supabase serves as the core Backend-as-a-Service (BaaS) infrastructure, providing a hosted PostgreSQL relational database for system state, metadata, and the audit log; Supabase Storage for large object files (the Document Vault); and Supabase Auth integrated natively with Row Level Security (RLS) policies.

Metadata Tagging and Phase-Aware Storage: When an applicant uploads a document, the physical file is stored exactly once in a dedicated Supabase Storage bucket under a unique, secure path. A Documents table records the file's metadata (e.g., \texttt{ApplicantID}, \texttt{DocumentType}, \texttt{Phase}, \texttt{UploadTimestamp}, \texttt{Verification\_Status}, and \texttt{Storage\_Path}), and the Applications table links to it by foreign key. A corrected file uploaded during a Conditional Revert is stored as a new version linked to the same requirement, preserving the history of what was submitted and when.

Document Verification Workflow: Every uploaded document, including proof-of-income documents such as BIR Form 2316 or a sworn statement, is automatically assigned a verification status of Pending. A University Scholarship Office staff member reviews the document through the Committee Portal and manually updates its status to Verified, or sets the application to Conditional Revert if the document is missing or illegible. Applications with Pending or unresolved documents cannot be endorsed to interview scheduling. All verification is performed internally by the University Scholarship Office; the platform integrates with no external verification service.

\subsection{Role-Based Access Control (RBAC) and Data Privacy Compliance}
\label{sec:rbac}

Because the ScholarPath AdDU platform processes highly sensitive economic and academic data, its architecture is governed by the Philippine Data Privacy Act of 2012 (RA 10173). Encryption and Row Level Security (RLS): All data in transit is secured using Transport Layer Security (TLS). Access to sensitive data at rest, in both database tables and Supabase Storage buckets, is restricted using RLS policies, so access controls are enforced directly at the database level regardless of how the frontend is accessed. RBAC Permissions Matrix: To enforce the principle of least privilege, users are assigned distinct privileges based on their role within the university, mapped directly to database RLS policies (Table~\ref{tab:rbac}).

\begin{table}[htbp]
  \centering
  \small
  \caption{ScholarPath AdDU Role-Based Access Control (RBAC) Matrix}
  \label{tab:rbac}
  \begin{tabular}{@{}p{3.2cm}p{5.2cm}p{6.0cm}@{}}
    \toprule
    User Role & Data Access Privileges & Core System Actions \\
    \midrule
    Applicant / Scholar & Read/write, self only & Register, choose a track and up to two programs, upload Phase~1 and Phase~2 documents, respond to Conditional Revert notices, view status and Academic Trajectory Path \\
    University Scholarship Office staff & Read/write, all applications & Verify documents, set Conditional Revert, confirm Lapsed disapprovals, record Layer~1 outcomes, endorse applications to interview, broadcast announcements \\
    Interview Panel & Read, assigned applications only; write, own scores & Record interview scores and recommendations \\
    School Scholarship Subcommittee & Read, own school's partition; write, own decisions & Filter and sort applicants, record selection decisions, confirm waitlist reallocations, send the accepted list to the Office of Admission \\
    Office of Admission & Read, accepted lists & Release results, which triggers applicant notifications \\
    System Administrator & Configuration only & Configure cycle windows, rule sets, and roles; no authority over applicant status \\
    \bottomrule
  \end{tabular}
\end{table}

\subsection{Automated Notification Subsystem}
\label{sec:notifications}

The absence of proactive alerts frequently results in missed submission windows and uncorrected files. ScholarPath AdDU integrates an automated, multi-channel notification subsystem that delivers deadline reminders, correction notices, and status updates via SMS and email. Third-Party API Integration: The SendGrid API formats and dispatches detailed email notifications, and the Twilio SMS API pushes brief, time-sensitive text messages to the applicant's registered mobile device. These services deliver messages only; they take no part in verification or decisions. Automated Trigger Logic and Edge Functions: The notification engine is driven by scheduled tasks (Cron jobs) and Supabase Edge Functions evaluating three conditions:

\begin{enumerate}
  \item Deadline Proximity: A scheduled worker queries the database daily. If the current date is within a predefined threshold (e.g., three days) of the close of the Phase~1 or Phase~2 window and the applicant has outstanding documents, an Edge Function sends a reminder via Twilio.
  \item Conditional Revert Clock: When staff set an application to Conditional Revert, an Edge Function sends the correction notice listing each missing or illegible file, followed by reminders within the five-day window. When the window elapses, the Lapsed flag triggers an alert to University Scholarship Office staff rather than a message to the applicant.
  \item Status Mutation: When a human actor commits a status change, a database webhook triggers an Edge Function. Selection outcomes are withheld from applicants until the Office of Admission records the release of results, after which the award, waitlist, or disapproval notice is sent via SendGrid. A waitlisted candidate is notified of a reallocated slot only after the subcommittee confirms the reallocation. % TODO(OPEN-8): notification timing after the Office of Admission release.
\end{enumerate}

\section{Testing \& Evaluation Procedures: ISO/IEC 25010 Task-Based Usability Protocol}
\label{sec:testing}

To assess the functional suitability and usability of the ScholarPath AdDU system, the study employs a mixed-method evaluation approach. This methodology combines task-based observational testing with a standardized System Usability Scale (SUS) questionnaire, administered to a purposive sample of approximately 10 to 15 end-users. The sample will consist of current or recent scholarship applicants, University Scholarship Office staff, and School Scholarship Subcommittee members, to guarantee coverage of the system's multi-role interfaces.

\subsection{Step-by-Step Testing Procedure}
\label{sec:procedure}

The evaluation will be conducted in a controlled environment following a standardized script to prevent researcher bias. The procedure is outlined as follows:

\begin{enumerate}
  \item Briefing and Consent: Participants are introduced to the system's purpose. The facilitator explicitly states that the system is being tested, not the user, and secures informed consent.
  \item Task Execution: Participants complete the tasks for their role. The facilitator observes silently without offering assistance unless a critical failure occurs.
  \begin{itemize}
    \item Task 1: Applicant Registration (applicants): The user creates an account, completes the application form, and selects one internal track and up to two degree programs.
    \item Task 2: Two-Stage Document Submission (applicants): The user uploads the Phase~1 documents, then the Phase~2 documents, and responds to a Conditional Revert notice by uploading a corrected file.
    \item Task 3: Verification and Conditional Revert (University Scholarship Office staff): The user verifies an application's documents, sets a Conditional Revert on an application with a missing file, and confirms the disapproval of a Lapsed application.
    \item Task 4: Committee Review and Reallocation (School Scholarship Subcommittee members): The user filters the school's partition, records a selection decision, and confirms a proposed waitlist reallocation after an awardee declines.
  \end{itemize}
  \item Observational Data Collection: During task execution, the facilitator records the usability metrics detailed in Section~\ref{sec:analysis}.
  \item Post-Session Instruments: Immediately after completing the tasks, participants complete the standardized 10-item System Usability Scale (SUS) questionnaire first, capturing their holistic impression before any comparative framing is introduced. Applicant participants then complete the Post-Prototype Perceived Effort Survey (Appendix~\ref{app:post-survey}); staff and subcommittee participants complete the Committee Portal Usefulness Survey (Appendix~\ref{app:committee-survey}). These instruments are not to be conflated: the SUS measures general interface usability, while the role-specific surveys generate the comparative and perception data used to answer Research Questions 1 through 3.
  \item Debriefing: A brief, semi-structured exit interview is conducted to gather qualitative feedback on any frustrations or difficulties encountered.
\end{enumerate}

\subsection{Data Analysis and Statistical Tools}
\label{sec:analysis}

To objectively quantify the system's performance and user satisfaction, the following metrics and statistical tools will be utilized:

\begin{itemize}
  \item Quantitative Observational Metrics:
  \begin{itemize}
    \item Time-on-Task: Measured in seconds, this records the duration a user takes to successfully complete each designated task.
    \item Error Rates: Calculated as the number of misclicks, failed uploads, or incorrect navigation paths taken per task, as well as the overall percentage of users who fail to complete a task independently.
    \item Notification Delivery Success Rate: Measured during Tasks 3 and 4, this metric records the percentage of triggered notification events (the Conditional Revert notice sent when a staff participant sets the status, and the outcome notice sent after a simulated Office of Admission release) that result in a successfully dispatched notification to the applicant's registered contact channel within the testing session. The facilitator confirms delivery on the designated applicant-side test device. The metric is computed as: (Number of Successfully Delivered Notifications / Total Triggered Notification Events) $\times$ 100. A delivery success rate of 100\% is the target threshold, given the controlled, low-concurrency testing environment.
    \item Trigger-to-Delivery Latency: Recorded in seconds, this measures the elapsed time between the moment a staff participant commits a status change (the event trigger) and the moment the corresponding notification is confirmed as received on the applicant-side test device. Mean latency across all triggered events in the session is computed and reported, operationalizing the performance efficiency of the notification subsystem as distinct from overall system response time.
  \end{itemize}
  \item System Usability Scale (SUS) Analysis:
  \begin{itemize}
    \item The SUS provides a validated instrument for measuring participants' subjective perceived usability of the platform. Responses are gathered on a standard 5-point Likert scale across 10 items, with alternating positive and negative phrasing to reduce response bias. Raw SUS scores, ranging from 0 to 100, will be analyzed using descriptive statistics. The Mean will be calculated to determine the average perceived usability score across the purposive sample. To contextualize the resulting mean SUS score, the study will apply the adjective rating scale validated by Bangor et al. \cite{ref13}: scores of 85 and above are classified as Excellent, 71--84 as Good, 50--70 as OK, and below 50 as indicative of usability failure requiring significant redesign. A target mean SUS score of 71 or higher (the Good threshold) is set as the minimum benchmark for the prototype to be considered functionally acceptable for its intended user base. The Standard Deviation will measure the variance in scores, indicating the degree of consensus among participants. The SUS addresses Research Question 4 specifically and is analytically distinct from the comparative instruments described below.
  \end{itemize}
  \item Functional Correctness Evaluation (Lifecycle Test Case Matrix):
  \begin{itemize}
    \item To assess whether the application status state machine (Listing~\ref{lst:lifecycle}) executes the Application Lifecycle Path correctly and never changes an applicant's status without a recorded human actor, corresponding to the Functional Correctness sub-characteristic of ISO/IEC 25010 (Table~\ref{tab:iso25010}), the study employs a structured test case matrix administered separately from the user-facing usability session (Table~\ref{tab:lifecycle-tests}). For each case, the expected system output is documented prior to test execution; the actual output is then recorded and compared. A match constitutes a pass; any discrepancy constitutes a failure and is logged for root cause analysis. The Functional Correctness score is computed as: (Number of Passing Test Cases / Total Test Cases) $\times$ 100. A pass rate of 100\% on the human-confirmation cases (TC-05, TC-07, TC-12, TC-13, and TC-16) is set as the non-negotiable threshold, because a status change without a recorded human actor is the highest-risk failure mode of the platform.
  \end{itemize}
  \item Pre/Post Comparative Survey Analysis:
  \begin{itemize}
    \item The Pre-Study Student Problem Validation Survey (Appendix~\ref{app:pre-survey}) and the Post-Prototype Perceived Effort Survey (Appendix~\ref{app:post-survey}) together constitute the comparative data collection instruments for Research Questions 1 and 3. Because both instruments use structurally parallel items with identical scale anchors, their data is analyzed through direct item-level comparisons using descriptive statistics.
    \item For Component 1 of both instruments (the ordinal time-estimation scale), the distribution of responses across the five ordinal ranges will be tabulated for both conditions. The modal response category and the direction of shift will be reported to address RQ1.
    \item For Component 2 of both instruments (the six-item perceived effort Likert grid), the Mean score per sub-task item will be calculated for both conditions and presented in a side-by-side comparison table. A decrease in mean score from pre- to post-prototype indicates a perceived reduction in effort for that application stage. Sub-task items 1 through 5 inform RQ1, and sub-task item 6 (deadline and correction tracking) informs RQ3.
    \item For Component 3 of both instruments, frequency distributions will be used for the categorical items concerning incomplete-document history and missed-deadline history (RQ3). The proportion of participants who reported an application delayed or rejected because of a missing or illegible document will be compared against the proportion who indicated that the Conditional Revert notice would let them correct such a file in time. Similarly, the proportion who reported missing deadlines will be compared against post-prototype responses on the notification subsystem's perceived effectiveness.
    \item For Component 4 of the post-prototype instrument (the overall comparative perception item), a frequency distribution of responses across the five-point scale will be reported as a summary indicator of the system's perceived overall impact.
  \end{itemize}
  \item Committee Portal Usefulness Analysis: Responses to the Committee Portal Usefulness Survey (Appendix~\ref{app:committee-survey}) will be summarized as item means and standard deviations to address RQ2(b).
  \item Qualitative Data Analysis (Thematic Analysis): The qualitative feedback gathered during the post-task semi-structured exit interviews will be processed using thematic analysis. The researchers will transcribe participant comments and categorize them into recurring themes (e.g., ``Navigation Friction,'' ``Terminology Confusion,'' or ``Positive Interface Feedback''). These themes will provide contextual explanations for the quantitative error rates and guide the final iterative refinements of the system's user interface prior to deployment.
\end{itemize}

\begin{longtable}{@{}p{1.2cm}p{6.6cm}p{6.6cm}@{}}
  \caption{Lifecycle Test Case Matrix}\label{tab:lifecycle-tests}\\
  \toprule
  ID & Test Profile & Expected System Output \\
  \midrule
  \endfirsthead
  \toprule
  ID & Test Profile & Expected System Output \\
  \midrule
  \endhead
  \bottomrule
  \endfoot
  TC-01 & High school average exactly 85\%, entrance exam ``ACCEPTED'', all Phase~1 documents & No baseline flag; application proceeds to Phase~2 \\
  TC-02 & High school average 84\% & Registration not blocked; application enters the staff review queue with a below-baseline flag; no status change until staff act \\
  TC-03 & Entrance exam result without an ``ACCEPTED'' remark & Baseline flag raised; application enters the staff review queue \\
  TC-04 & Jubilee applicant without a Rank~1 or Rank~2 certificate & Missing-document flag on the Jubilee requirement; application cannot be endorsed \\
  TC-05 & Phase~2 document missing; staff set Conditional Revert & Status becomes Conditional Revert; notice lists the missing file; audit log records the staff actor \\
  TC-06 & Corrected file uploaded on day~5 of the Revert window & Application returns to verification; no Lapsed flag \\
  TC-07 & No correction by day~6 & Status becomes Lapsed and staff are alerted; status does not become Disapproved until a staff member confirms \\
  TC-08 & SON applicant with a conditional entrance exam score & SON rule set flags the applicant; no selection is recorded by the system \\
  TC-09 & SEA applicant below the entrance exam cutoff & SEA rule set flags the applicant as not meeting the program minimum; decision remains with the subcommittee \\
  TC-10 & Applicant whose two program choices belong to different schools & Application appears in both schools' partitions \\
  TC-11 & Subcommittee records an award before the Office of Admission release & Status is recorded; no notification reaches the applicant until the release event \\
  TC-12 & Awardee declines the offer & Slot freed; subcommittee alerted with a proposed waitlisted candidate; no award to the candidate until the subcommittee confirms \\
  TC-13 & Attempt by a system process to set Disapproved, Waitlisted, or Awarded & Transition rejected by the human-actor guard \\
  TC-14 & Inter-track realignment from Working Scholars to GIA & Track changed by an evaluator; prior track retained in history and audit log \\
  TC-15 & Subcommittee member attempts to view another school's partition & Access denied by RLS \\
  TC-16 & Any completed test case & Every status change in the audit log carries an actor, timestamp, prior status, and new status \\
\end{longtable}

\subsection{Pre-Study, Post-Prototype, and Committee Survey Instruments}
\label{sec:instruments}

To operationalize the within-subjects comparative design described in Section~\ref{sec:research-design}, this study employs two structurally paired applicant instruments administered at distinct phases of the research process, together with a single-administration staff instrument. The Pre-Study Student Problem Validation Survey (Appendix~\ref{app:pre-survey}) is administered prior to system exposure to establish baseline data on the perceived burden of the existing paper-based process. The Post-Prototype Perceived Effort Survey (Appendix~\ref{app:post-survey}) is administered immediately following the task-based usability testing session using an identical scale structure. Together, these instruments generate the comparative data required to address Research Questions 1 and 3. Because each participant's pre-study responses serve as their own baseline, any observed reduction in perceived effort can be attributed to the system's features rather than to differences in prior scholarship experience. Both instruments use scaled and categorical response formats rather than open recall of specific durations, acknowledging that the manual application process unfolds over days or weeks and is not conducive to precise time measurement.

Pre-Study Student Problem Validation Survey (Appendix~\ref{app:pre-survey}). The pre-study instrument is administered to a purposive sample of current or recent AdDU scholarship applicants prior to any system exposure. It is organized into three components. Component 1: Perceived Time-on-Process asks respondents to estimate the total duration of their end-to-end manual application experience, from reading the scholarship announcement through physically submitting the application form and supporting documents, using five ordinal ranges: less than 1 day, 1--3 days, 4--7 days, 1--2 weeks, and more than 2 weeks. This item establishes the time-burden baseline for RQ1. Post-prototype time estimates are derived from a controlled testing session and may not reflect the duration of the process in real-world deployment, where document readiness, network conditions, and platform familiarity may produce different results. Component 2: Perceived Effort Scale applies a five-point Likert format (1 = Very Easy, 5 = Very Difficult) to six application sub-tasks: (1) understanding which scholarship track to apply for and what it requires; (2) determining whether they met the baseline requirements, such as the high school average, entrance exam result, rank certificate, and income documents, before submitting; (3) gathering, printing, and preparing the required physical documents; (4) physically submitting the completed application package; (5) tracking the status of the application after submission; and (6) keeping track of submission deadlines and requests for missing documents. Sub-tasks 1 through 5 establish the effort baseline for RQ1, and sub-task 6 establishes the deadline and correction baseline for RQ3. Component 3: Document and Deadline History uses categorical items to establish behavioral baselines for RQ3. Respondents are asked whether an application of theirs was ever delayed or rejected because of a missing or illegible document, and how often; and whether they ever missed a scholarship application deadline, and for what primary reason.

Post-Prototype Perceived Effort Survey (Appendix~\ref{app:post-survey}). The post-prototype instrument is administered to the same participants after the SUS questionnaire. Its structural mirroring of the pre-study instrument is deliberate: because both instruments use identical Likert anchors, ordinal ranges, and sub-task categories, their corresponding items can be compared directly. Component 1 presents the same five ordinal time-estimation ranges. Component 2 presents the same six sub-tasks, reworded to reflect ScholarPath AdDU's features: track and program selection at registration; understanding the Phase~1 Pre-qualification result; uploading documents through the Document Vault; submitting the application digitally; tracking status through the applicant dashboard; and staying aware of deadlines and correction requests through the notification subsystem. Component 3: System Feature Effectiveness asks whether the Conditional Revert notice would help participants correct a missing or illegible document before their application lapses, and how far the automated notifications would help them remain aware of upcoming deadlines. Component 4: Overall Comparative Perception asks participants to characterize the total effort required by ScholarPath AdDU relative to the paper-based process on a five-point scale. An optional open-ended item closes the instrument.

Committee Portal Usefulness Survey (Appendix~\ref{app:committee-survey}). University Scholarship Office staff and School Scholarship Subcommittee members who complete Tasks 3 and 4 rate, on a five-point agreement scale, whether the Committee Portal presents the information they need to decide, whether the school partition and filters reduce manual collation, whether the audit log makes decisions traceable, and whether they retain full control over every status decision. These items address RQ2(b).

\section{Ethical Considerations}
\label{sec:ethics}

Because the usability testing of ScholarPath AdDU involves human participants and the simulated input of sensitive socioeconomic data, the evaluation phase is strictly governed by established research ethics and the Philippine Data Privacy Act of 2012 (RA 10173). The following ethical protocols will be implemented throughout the testing phase:

\begin{itemize}
  \item Informed Consent and Voluntary Participation: Prior to any system interaction, all participants will be provided with an informed consent form detailing the purpose of the study, the specific tasks required, and the estimated time commitment. Because scholarship applicants may be incoming first-year students under 18 years of age, participants who are minors will additionally require the written consent of a parent or guardian together with their own assent. Participation is strictly voluntary, and participants maintain the right to withdraw from the testing session at any point without penalty or need for justification. % TODO(L-4): confirm minor-consent requirements with the ethics adviser.
  \item Data Anonymization and Privacy: During task execution, participants will be explicitly instructed to use provided mock data or pseudonymized profiles rather than their actual household income, grades, or documents. This ensures that no genuine sensitive information is exposed or recorded during the evaluation.
  \item Confidentiality and Secure Data Disposal: All observational metrics, survey responses (Appendices~\ref{app:pre-survey} to~\ref{app:committee-survey}), and SUS questionnaire responses will be kept strictly confidential. Data will be aggregated so that no single metric or quote can be traced back to an individual participant. Upon the successful defense and final submission of the capstone project, all test accounts, mock databases, and physical consent forms utilized during the evaluation phase will be permanently and securely deleted.
\end{itemize}
